% !TeX root = ../drafts/07-instruction-tables.tex
\section{The instruction tables}\label{sec:tables}

Each instruction has its own table.

Write $\rdf{\sep,\addr,\cnt,\val_0,\val_1, \dots}$ for the pair \pull\ $\tup{\sep,\addr,\cnt,\val_0,\val_1, \dots }$, \push\ $\tup{\sep,\addr,\gen\cdot\cnt,\val_0,\val_1, \dots}$: every flush below but the state's.

\subsection{\texttt{XOR}}\label{sec:tab-xor}

\texttt{XOR}\ $[o_A,o_B,o_C]$ asserts $\loc{o_C}=\loc{o_A}+\loc{o_B}$, the $192$-bit field sum (in $\E$).
\begin{itemize}
\item \textbf{Columns:} $\pc,\fp$; operands $o_A,o_B,o_C$; the two read values $(v_{A,i})_{i\in\{0,1,2\}}$ and $(v_{B,i})_{i\in\{0,1,2\}}$; memory counts $r_A,r_B,r_C$; bytecode count $r_{\mathrm{bc}}$.
\item \textbf{Constraints:} none.
\item \textbf{Flushes:}
  \begin{itemize}
  \item \textbf{State:} \pull\ $\tup{\dsep{ST},\pc,\fp}$, \push\ $\tup{\dsep{ST},\gen\cdot\pc,\fp}$.
  \item \textbf{Bytecode:} $\rdf{\dsep{BC},\pc,r_{\mathrm{bc}},\opc{XOR},o_A,o_B,o_C,0,0}$.
  \item \textbf{Memory:}
    \begin{align*}
      &\rdf{\dsep{MEM},\fp\cdot o_X,r_X,v_{X,0},v_{X,1},v_{X,2}}\quad(X\in\{A,B\}),\\
      &\rdf{\dsep{MEM},\fp\cdot o_C,r_C,v_{A,0}+v_{B,0},v_{A,1}+v_{B,1},v_{A,2}+v_{B,2}}.
    \end{align*}
  \end{itemize}
\end{itemize}

\subsection{\texttt{MUL\_NATIVE}}\label{sec:tab-mul}

\texttt{MUL\_NATIVE}\ $[o_A,o_B,o_C]$ asserts $\loc{o_C}=\loc{o_A} \cdot \loc{o_B}$, the $192$-bit field multiplication (in $\E$).
\begin{itemize}
\item \textbf{Columns:} $\pc,\fp$; operands $o_A,o_B,o_C$; the two read values $(v_{A,i})_{i\in\{0,1,2\}}$ and $(v_{B,i})_{i\in\{0,1,2\}}$; memory counts $r_A,r_B,r_C$; bytecode count $r_{\mathrm{bc}}$.
\item \textbf{Constraints:} none.
\item \textbf{Flushes:}
  \begin{itemize}
  \item \textbf{State:} \pull\ $\tup{\dsep{ST},\pc,\fp}$, \push\ $\tup{\dsep{ST},\gen\cdot\pc,\fp}$.
  \item \textbf{Bytecode:} $\rdf{\dsep{BC},\pc,r_{\mathrm{bc}},\opc{MUL},o_A,o_B,o_C,0,0}$.
  \item \textbf{Memory} (the destination carries $v_Av_B$ in $\E$, of degree two). With $y^3=y+1$, put
    \[
    p_0=v_{A,0}v_{B,0},\quad
    p_1=v_{A,0}v_{B,1}+v_{A,1}v_{B,0},\quad
    p_2=v_{A,0}v_{B,2}+v_{A,1}v_{B,1}+v_{A,2}v_{B,0},
    \]
    \[
    p_3=v_{A,1}v_{B,2}+v_{A,2}v_{B,1},\qquad
    p_4=v_{A,2}v_{B,2},
    \]
    twelve products over the nine pairs $v_{A,i}v_{B,j}$:
    \begin{align*}
      &\rdf{\dsep{MEM},\fp\cdot o_X,r_X,v_{X,0},v_{X,1},v_{X,2}}\quad(X\in\{A,B\}),\\
      &\rdf{\dsep{MEM},\fp\cdot o_C,r_C,p_0+p_3,p_1+p_3+p_4,p_2+p_4}.
    \end{align*}
  \end{itemize}
\end{itemize}

\subsection{\texttt{SET\_CONSTANT}}\label{sec:tab-set}

\texttt{SET\_CONSTANT}\ $[o,k_0,k_1,k_2]$ asserts $\loc{o}=k_0+k_1y+k_2y^{2}$.
\begin{itemize}
\item \textbf{Columns:} $\pc,\fp$; operand $o$; immediate $(k_i)_{i\in\{0,1,2\}}$; memory count $r$; bytecode count $r_{\mathrm{bc}}$.
\item \textbf{Constraints:} none.
\item \textbf{Flushes:}
  \begin{itemize}
  \item \textbf{State:} \pull\ $\tup{\dsep{ST},\pc,\fp}$, \push\ $\tup{\dsep{ST},\gen\cdot\pc,\fp}$.
  \item \textbf{Bytecode:} $\rdf{\dsep{BC},\pc,r_{\mathrm{bc}},\opc{SET},o,k_0,k_1,k_2,0}$.
  \item \textbf{Memory:} $\rdf{\dsep{MEM},\fp\cdot o,r,k_0,k_1,k_2}$.
  \end{itemize}
\end{itemize}

\subsection{\texttt{DEREF}}\label{sec:tab-deref}

\texttt{DEREF}\ $[o_1,o_2,o_3;\,\dmode]$ reads a pointer and asserts the cell it points to. The row reads three cells:
\begin{itemize}
\item the \emph{pointer} $p$, at $\fp\cdot o_1$;
\item the \emph{target} $v_2$, at the dereferenced address $p\cdot o_2$;
\item the \emph{local} value $v_3$, at $\fp\cdot o_3$.
\end{itemize}
The store asserts that $v_2$ equals the source chosen by the mode $\dmode$ (\S\ref{sec:vm}): the local cell $v_3$, the return address $\gen^{2}\cdot\pc$, or the frame pointer $\fp$. The mode is two boolean flags $(f_{\pc},f_{\fp})$: $(0,0)$ for $\texttt{deref\_cell}[o_3]$, $(1,0)$ for $\texttt{deref\_pc}$, $(0,1)$ for $\texttt{deref\_fp}$.
\begin{itemize}
\item \textbf{Columns:} $\pc,\fp$; operands $o_1,o_2,o_3$; flags $f_{\pc},f_{\fp}$; pointer $p$ (a single $\K$-limb); the local value $(v_{3,i})_{i\in\{0,1,2\}}$; memory counts $r_1,r_2,r_3$; bytecode count $r_{\mathrm{bc}}$.
\item \textbf{Constraints:} none.
\item \textbf{Flushes:}
  \begin{itemize}
  \item \textbf{State:} \pull\ $\tup{\dsep{ST},\pc,\fp}$, \push\ $\tup{\dsep{ST},\gen\cdot\pc,\fp}$.
  \item \textbf{Bytecode:} $\rdf{\dsep{BC},\pc,r_{\mathrm{bc}},\opc{DRF},o_1,o_2,o_3,f_{\pc},f_{\fp}}$.
  \item \textbf{Memory}, the target giving $v_3$, $\gen^{2}\cdot\pc$ or $\fp$ at the three flag settings:
    \begin{align*}
      &\rdf{\dsep{MEM},\fp\cdot o_1,r_1,p,0,0},\\
      &\rdf{\dsep{MEM},\fp\cdot o_3,r_3,v_{3,0},v_{3,1},v_{3,2}},\\
      &\rdf{\dsep{MEM},p\cdot o_2,r_2,\bar f\,v_{3,0}+f_{\pc}\,(\gen^{2}\cdot\pc)+f_{\fp}\,\fp,\ \bar f\,v_{3,1},\ \bar f\,v_{3,2}},
      \qquad \bar f=1+f_{\pc}+f_{\fp}.
    \end{align*}
  \end{itemize}
\end{itemize}

\subsection{\texttt{JUMP}}\label{sec:tab-jump}

\texttt{JUMP}\ $[o_c,o_d,o_f]$ performs a conditional jump on whether the condition $v_{\mathrm{cond}}=\loc{o_c}$ is nonzero. It transfers to $\pc\gets v_{\pc} = \loc{o_d}$, $\fp\gets v_{\fp} = \loc{o_f}$ when $v_{\mathrm{cond}}\neq0$, and defaults to $\pc\gets\gen\cdot\pc$, $\fp\gets\fp$ otherwise. $v_{\mathrm{cond}}, v_{\fp}, v_{\pc}$ must be $\K$-valued. The branch is taken according to a boolean indicator column $b=[\,v_{\mathrm{cond}}\neq0\,]$, certified by a prover-supplied inverse $w$ and 2 constraints:

\begin{itemize}
\item \textbf{Columns:} $\pc,\fp$; operands $o_c,o_d,o_f$; the condition $v_{\mathrm{cond}}$, the destination $v_{\pc}$ and the frame $v_{\fp}$; inverse $w$ (honest prover sets $w = v_{\mathrm{cond}}^{-1}$ if $v_{\mathrm{cond}}\neq0$, otherwise $w=0$); boolean indicator $b$; memory counts $r_c,r_d,r_f$; bytecode count $r_{\mathrm{bc}}$.
\item \textbf{Constraints:}  $b=v_{\mathrm{cond}}\cdot w$ and $v_{\mathrm{cond}}\,(b+1)=0$
\item \textbf{Flushes:}
  \begin{itemize}
  \item \textbf{State:} \pull\ $\tup{\dsep{ST},\pc,\fp}$, \push\ $\tup{\dsep{ST},b\,v_{\pc}+b\,(\gen\cdot\pc)+\gen\cdot\pc,\;b\,v_{\fp}+b\,\fp+\fp}$.
  \item \textbf{Bytecode:} $\rdf{\dsep{BC},\pc,r_{\mathrm{bc}},\opc{JMP},o_c,o_d,o_f,0,0}$.
  \item \textbf{Memory:}
    \begin{align*}
      &\rdf{\dsep{MEM},\fp\cdot o_c,r_c,v_{\mathrm{cond}},0,0},\\
      &\rdf{\dsep{MEM},\fp\cdot o_d,r_d,v_{\pc},0,0},\\
      &\rdf{\dsep{MEM},\fp\cdot o_f,r_f,v_{\fp},0,0}.
    \end{align*}
  \end{itemize}
\end{itemize}

\subsection{\texttt{BLAKE2S}}\label{sec:tab-blake2s}

\texttt{BLAKE2S}\ $[o_{m_0},o_{m_1},o_{m_2},o_{m_3},o_{\mathit{cv}},o_{\mathit{out}},\md_0,\md_1]$ compresses a $64$-byte message block under a $32$-byte chaining value into a $32$-byte result, with BLAKE2s flags and byte counter hardcoded in the metadata $(\md_0,\md_1)$. Each $128$-bit half occupies one cell, top limb zero: the four message chunks at $\fp\cdot o_{m_0},\dots,\fp\cdot o_{m_3}$, addressed independently, the chaining value at the consecutive pair $\fp\cdot o_{\mathit{cv}},\gen\cdot\fp\cdot o_{\mathit{cv}}$, and the result at $\fp\cdot o_{\mathit{out}},\gen\cdot\fp\cdot o_{\mathit{out}}$.

Flock~\cite{Flock26} proves the compression relation with a batched R1CS over a Boolean witness (Annex~\ref{annex:flock}). The witness is packed into $\qflock$ via Ring-Switching (Annex~\ref{annex:rs}), with $64$ bits per $\K$-element.

The message, result, chaining value and metadata are all committed in $\qflock$: the two low limbs $z_0,z_1$ of each of the eight cells $z\in\{m_0,\dots,m_3,\mathit{cv}_0,\mathit{cv}_1,\mathit{out}_0,\mathit{out}_1\}$, plus $\md_0,\md_1$, eighteen in all. Linking these values to the memory interaction implicitly forces the top limb of every memory cell accessed by BLAKE2s to be zero. The generalized R1CS matrices take the chaining value, byte counter, and the two finalization flags as free input: the bytecode interactions via the bus bind them.
\begin{itemize}
\item \textbf{Columns:} $\pc,\fp$; operands $o_{m_0},o_{m_1},o_{m_2},o_{m_3},o_{\mathit{cv}},o_{\mathit{out}}$; one memory count $r_z$ per cell $z$; bytecode count $r_{\mathrm{bc}}$. The eighteen limbs are committed in $\qflock$, no duplicate columns needed. The eight addresses are the products $\gen^{k}\fp\cdot o$, $k\in\{0,1\}$.
\item \textbf{Constraints:} none. The compression relation is handled by Flock.
\item \textbf{Flushes:}
  \begin{itemize}
  \item \textbf{State:} \pull\ $\tup{\dsep{ST},\pc,\fp}$, \push\ $\tup{\dsep{ST},\gen\cdot\pc,\fp}$.
  \item \textbf{Bytecode:} $\rdf{\dsep{BC},\pc,r_{\mathrm{bc}},\opc{B2S},o_{m_0},o_{m_1},o_{m_2},o_{m_3},o_{\mathit{cv}},o_{\mathit{out}},\md_0,\md_1}$.
  \item \textbf{Memory:}
    \begin{align*}
      &\rdf{\dsep{MEM},\fp\cdot o_{m_i},r_{m_i},m_{i,0},m_{i,1},0}\quad(i=0,\dots,3),\\
      &\rdf{\dsep{MEM},\fp\cdot o_{\mathit{cv}},r_{\mathit{cv}_0},\mathit{cv}_{0,0},\mathit{cv}_{0,1},0},\\
      &\rdf{\dsep{MEM},\gen\cdot\fp\cdot o_{\mathit{cv}},r_{\mathit{cv}_1},\mathit{cv}_{1,0},\mathit{cv}_{1,1},0},\\
      &\rdf{\dsep{MEM},\fp\cdot o_{\mathit{out}},r_{\mathit{out}_0},\mathit{out}_{0,0},\mathit{out}_{0,1},0},\\
      &\rdf{\dsep{MEM},\gen\cdot\fp\cdot o_{\mathit{out}},r_{\mathit{out}_1},\mathit{out}_{1,0},\mathit{out}_{1,1},0}.
    \end{align*}
  \end{itemize}
\end{itemize}
